\section{Introduction}
\label{sec:intro}

De novo motif discovery --- the identification of recurring DNA sequence patterns from collections of functional genomic regions --- remains a central problem in computational biology. Given a set of DNA sequences enriched for transcription factor (TF) binding sites (typically ChIP-seq peaks), the goal is to discover position weight matrices (PWMs) that represent the binding specificity of the targeted TF.

Several algorithms have been developed over three decades of research. MEME \citep{bailey1994meme} pioneered expectation-maximization (EM) approaches and remains a gold standard for accuracy, but its computational cost makes it impractical for large datasets. STREME \citep{bailey2021streme}, the successor to DREME \citep{bailey2011dreme}, uses suffix-tree enumeration for faster discovery and is the current default in the MEME Suite. HOMER \citep{heinz2010homer} uses a distinct cumulative approach and is widely used in ChIP-seq analysis pipelines despite its slower runtime.

Despite this maturity, there remains a tension between accuracy and speed: the most accurate tools (MEME, HOMER) are slow, while faster tools (STREME) leave accuracy on the table. For large-scale analyses --- such as the ENCODE Consortium's \citep{encode2012} 500+ TF ChIP-seq datasets --- this trade-off imposes a practical bottleneck.

We present motif-discover, a tool that substantially narrows this gap. On the standard benchmark of dinucleotide-shuffled negatives, motif-discover significantly outperforms STREME in both accuracy and speed. On the more challenging genomic negatives benchmark, it matches MEME's accuracy --- the gold standard --- while running approximately 300$\times$ faster. The tool is distributed as a single statically-linked binary with no dependencies, enabling reproducible motif analysis on any x86-64 Linux system.
